\subsection{Attachment of a topological space}

Let X, B be topological spaces, let A be a subspace of B and let \(f: A \rightarrow X\) be a continuous map. Let Y denote the topological sum of the spaces \(Y_{1} = X\) and \(Y_{2} = B\) and identify \(Y_{1}\) and \(Y_{2}\) with subspaces of Y by the canonical injections. Let R be the equivalence relation in Y which is the finest for which the elements a of \(Y_{2}\) and \(f(a)\) of \(Y_{1}\) are equivalent, for every \(a \in A\) . The relation R is the equality relation in \(Y_{1}\) . Let \(x \in Y_{1}\) and let \(b \in Y_{2}\) ; we have x R b if and only if \(b \in A\) and \(f(b) = x\) . Let \(b, b' \in Y_{2}\) ; in order that one have b R \(b'\) , it is necessary and sufficient that \(b = b'\) , or that \(b, b' \in A\) and \(f(b) = f(b')\) .

DEFINITION 9. --- The quotient topological space Y/R is called the space obtained by attaching to the space X the space B along the map f. It is denoted \(X \cup_{f} B\) .

Denote by \(\alpha_{f}\colon\mathrm{X}\to\mathrm{X}\cup_{f}\mathrm{B}\) and \(\beta_{f}\colon\mathrm{B}\to\mathrm{X}\cup_{f}\mathrm{B}\) the restrictions to \(\mathrm{Y}_{1}\) and \(\mathrm{Y}_{2}\) of the canonical surjection of Y onto Y/R. We have \(\alpha_{f}\circ f=\beta_{f}\mid\mathrm{A}\) .

Remarks. --- 1) The map \(\alpha_{f}\) is injective. For every subset U of X, we have \(\bar{\alpha}_{f}^{-1}(\alpha_{f}(U)) = U\) and \(\bar{\beta}_{f}^{-1}(\alpha_{f}(U)) = f^{-1}(U)\) . If A is an open (resp. closed) subset of B, these relations imply that \(\alpha_{f}\) is an open (resp. closed) map from X into \(X \cup_{f} B\) .

Let U be an open subset of X and let V be an open set of B such that \(f^{-1}(U) = V \cap A\) . The union of the subsets U of \(Y_{1}\) and V of \(Y_{2}\) is an open saturated subset of Y; its image in \(X \cup_{f} B\) is therefore open and its trace on \(\alpha_{f}(X)\) is \(\alpha_{f}(U)\) . Consequently, the map \(\alpha_{f}\) defines a homeomorphism of X onto its image in \(X \cup_{f} B\) .

\begin{enumerate}
\def\labelenumi{\arabic{enumi})}
\setcounter{enumi}{1}
\tightlist
\item
  Let V be a subset of B. We have \(\bar{\alpha}_{f}^{-1}(\beta_{f}(V)) = f(V \cap A)\) and \(\bar{\beta}_{f}^{-1}(\beta_{f}(V)) = V \cup \bar{f}(f(V \cap A))\) . If A is closed in B and if the map f is closed, the map \(\beta_{f}\) is therefore closed. Similarly, if A is open in B and if the map f is open, the map \(\beta_{f}\) is open.
\end{enumerate}

In these two cases, the map \(\beta_{f}\) then induces a homeomorphism of \(B\setminus A\) onto its image.

\begin{enumerate}
\def\labelenumi{\arabic{enumi})}
\setcounter{enumi}{2}
\tightlist
\item
  Suppose that A is closed and that the map \(f\) is proper. We have just seen that the map \(\beta_f\) is closed. For every \(x \in X\) , \(\overline{\beta_f}(\alpha_f(x)) = \overline{f}(x)\) . It is therefore a quasi-compact subset of A (TG, I, p. 75, theorem 1), hence also of B. For every \(b \in B\) , \(\overline{\beta_f}(\beta_f(b))\) equal to \(\{b\}\) if \(b \in B - A\) , and \(\overline{f}(f(b))\) if \(b \in A\) ; in both cases, it is a quasi-compact subset of B. The fibers of the map \(\beta_{f}\) are therefore quasi-compact; consequently (loc. cit.), the map \(\beta_{f}\) is proper.
\end{enumerate}

The map \(\alpha_{f}\) is also proper (TG, I, p. 72, prop. 2). It follows that the canonical map of Y onto X \(\cup_{f}\) B is proper.

Examples. --- 1) Let X and Y be topological spaces and \(f\colon X\to Y\) a continuous map. The cylinder \(\operatorname{Cyl}(f)\) is none other than the space obtained by attaching to the space Y the space \(X\times I\) along the map \(f\circ \mathrm{pr}_1\) of \(X\times \{1\}\) in Y.

\begin{enumerate}
\def\labelenumi{\arabic{enumi})}
\setcounter{enumi}{1}
\tightlist
\item
  Let X be a topological space, let B be a topological space and let A be a subspace of B, let \(f: A \to X\) be a continuous map which induces a homeomorphism of A onto its image.
\end{enumerate}

Let us set \(\mathrm{A}' = f(\mathrm{A})\); let \(f'\) be the map from \(\mathrm{A}'\) into B which associates to every \(x \in \mathrm{A}'\) the unique preimage of \(x\) under \(f\). The map \(f' \circ f\) is continuous; since \(f\) is strict, the map \(f'\) is continuous. There exists a unique continuous map \(u\) from the space \(\mathrm{X} \cup_f \mathrm{B}\) to the space \(\mathrm{B} \cup_{f'} \mathrm{X}\) which maps \(\alpha_f(x)\) onto \(\beta_{f'}(x)\) for \(x \in \mathrm{X}\) and \(\beta_f(b)\) onto \(\alpha_{f'}(b)\) for \(b \in \mathrm{B}\); it is a homeomorphism which maps the subspace \(\alpha_f(\mathrm{X})\) onto the subspace \(\beta_{f'}(\mathrm{X})\).

PROPOSITION 8. --- Let Z be a topological space.

\begin{enumerate}
\def\labelenumi{\alph{enumi})}
\item
  Let \(u \colon X \to Z\) and \(v \colon B \to Z\) be continuous maps such that \(v(a) = u(f(a))\) for all \(a \in A\). There then exists a unique continuous map \(w \colon X \cup_f B \to Z\) such that \(w \circ \alpha_f = u\) and \(w \circ \beta_f = v\).
\item
  Let \(\sigma\colon\mathbf{X}\times\mathbf{I}\to\mathbf{Z}\) and \(\tau\colon\mathbf{B}\times\mathbf{I}\to\mathbf{Z}\) be homotopies. It is assumed that \(\tau\) is fixed on A and that \(\sigma(f(a),t)=\tau(a,t)\) for all \(a\in\mathbf{A}\) and every \(t\in\mathbf{I}\). There then exists a unique homotopy \(\eta\colon(\mathbf{X}\cup_{f}\mathbf{B})\times\mathbf{I}\to\mathbf{Z}\) such that \(\eta(\alpha_{f}(x),t)=\sigma(x,t)\) and \(\eta(\beta_{f}(b),t)=\tau(b,t)\) for \(x\in\mathbf{X}\), \(b\in\mathbf{B}\), and \(t\in\mathbf{I}\). This homotopy is fixed on \(\beta_{f}(\mathbf{A})\).
\end{enumerate}

The proposition follows immediately from the definition of a quotient space and from Prop. 10 (I, p. 19).

COROLLARY 1. --- Let \(B'\) be a topological space, let \(A'\) be a subspace of \(B'\), and let \(f': A' \to X\) be a continuous map. Let \(v: B \to B'\) be a continuous map such that \(v(A) \subset A'\) and \(f = f' \circ (v \mid A)\). Let \(w: X \cup_f B \to X \cup_{f'} B'\) be the unique continuous map such that \(w \circ \alpha_f = \alpha_{f'}\) and \(w \circ \beta_f = \beta_{f'} \circ v\). If \(v\) defines a homotopy equivalence of the pair \((B, A)\) onto the pair \((B', A')\), then \(w\) defines a homotopy equivalence of the pair \((X \cup_f B, \alpha_f(X))\) onto the pair \((X \cup_{f'} B', \alpha_{f'}(X))\).

Let \(v' \colon \mathrm{B}' \to \mathrm{B}\) be a continuous map, let \(\tau \colon \mathrm{B} \times \mathbf{I} \to \mathrm{B}\) be a homotopy fixed on A connecting \(\mathrm{Id}_{\mathrm{B}}\) to \(v' \circ v\), and let \(\tau' \colon \mathrm{B}' \times \mathbf{I} \to \mathrm{B}'\) be a homotopy fixed on \(\mathrm{A}'\) connecting \(\mathrm{Id}_{\mathrm{B}'}\) to \(v \circ v'\). For \(a' \in \mathrm{A}'\) and \(a = v'(a')\), we have \(a' = v(a)\) and \(\beta_f(v'(a')) = \beta_f(a) = \alpha_f(f(a)) = \alpha_f(f'(a'))\). By Proposition 8(a), there exists a unique map \(w' \colon X \cup_{f'} B' \to X \cup_f B\) such that \(w' \circ \alpha_{f'} = \alpha_{f'}\) and \(w' \circ \beta_{f'} = \beta_f \circ v'\). By Proposition 8(b), there exists a unique homotopy \(\eta \colon (\mathrm{X} \cup_f \mathrm{B}) \times \mathbf{I} \to X \cup_{f'} B'\) such that \(\eta(\alpha_f(x), t) = \alpha_{f'}(x)\) and \(\eta(\beta_f(b), t) = \beta_{f'}(\tau(b, t))\) for \(x \in X\), \(b \in B\), and \(t \in I\). There likewise exists a unique homotopy \(\eta' \colon (X \cup_{f'} B') \times I \to X \cup_f B\) such that \(\eta'(\alpha_{f'}(x), t) = \alpha_f(x)\) and \(\eta'(\beta_{f'}(b), t) = \beta_f(\tau'(b, t))\) for \(x \in X\), \(b \in B'\), and \(t \in I\). The map from \((X \cup_f B) \times I\) into \(X \cup_f B\) given by \((x, t) \mapsto \eta'(\eta(x, t), t)\) is then a homotopy fixed on \(\alpha_f(X)\) connecting the identity map of \(X \cup_f B\) to the map \(w' \circ w\). Similarly, the map from \((X \cup_{f'} B') \times I\) into \(X \cup_{f'} B'\) given by \((x, t) \mapsto \eta(\eta'(x, t), t)\) is a homotopy fixed on \(\alpha_{f'}(X)\) connecting the identity map of \(X \cup_{f'} B'\) to the map \(w \circ w'\). The corollary follows from this.

Example 3. --- Let i denote the canonical injection of A into B and take as space B' the mapping cylinder of i; denote \(\alpha_{i}\colon\mathrm{A}\times\mathbf{I}\to\mathrm{Cyl}(i)\) and \(\beta_{i}\colon\mathrm{B}\to\mathrm{Cyl}(i)\) the canonical maps and \(r_{i}\colon\mathrm{Cyl}(i)\to\mathrm{B}\) the canonical retraction of the cylinder \(\mathrm{Cyl}(i)\) on its base. Let \(A_{0}\) be the subspace \(\alpha_{i}(\mathrm{A}\times\{0\})\) of \(\mathrm{Cyl}(i)\) and denote by \(f_{0}\colon\mathrm{A}_{0}\to\mathrm{X}\) the map \(f\circ r_{i}\). For \(a\in A\), we have

\[
\beta_ {f} \circ r _ {i} (\alpha_ {i} (a, 0)) = \beta_ {f} (a) = \alpha_ {f} (f _ {0} (\alpha_ {i} (a, 0))) = \alpha_ {f} (f (a)).
\]

Let \(\eta\colon X\cup_{f_0}\mathrm{Cyl}(i)\to X\cup_fB\) be the unique continuous map such that \(\eta\circ\alpha_{f_0}=\alpha_f\) and \(\eta\circ\beta_{f_0}=\beta_f\circ r_i\). Suppose the pair (B, A) has the homotopy extension property. Then the map \(r_i\) defines a homotopy equivalence of the pair \((\mathrm{Cyl}(i),\mathrm{A}_0)\) onto the pair (B,A) (III, p. 243, th. 1). By Corollary 1, the map \(\eta\) then defines a homotopy equivalence of the pair \((X\cup_{f_0}\mathrm{Cyl}(i),\alpha_{f_0}(X))\) onto the pair \((X\cup_fB,\alpha_f(X))\). In particular, \(\eta\) is a homotopy equivalence.

COROLLARY 2. --- Let \(X'\) be a topological space, let \(u: X \to X'\) be a continuous map, and set \(f' = u \circ f\). Let \(w: X \cup_f B \to X' \cup_{f'} B\) be the unique continuous map such that \(w \circ \alpha_f = \alpha_{f'} \circ u\) and \(w \circ \beta_f = \beta_{f'}\). If \(u\) defines a homotopy equivalence of the pair \((X, f(A))\) onto the pair \((X', f'(A))\), then \(w\) defines a homotopy equivalence of the pair \((\mathrm{X}\cup_f\mathrm{B},\beta_f(\mathrm{B}))\) onto the pair \((\mathrm{X}'\cup_{f'}\mathrm{B},\beta_{f'}(\mathrm{B}))\).

The proof is analogous to that of Corollary 1.

PROPOSITION 9. --- Let X and B be topological spaces, let A be a subspace of B, and let f: A → X be a continuous map.

\begin{enumerate}
\def\labelenumi{\alph{enumi})}
\item
  If the pair (B, A) has the homotopy extension property, then so does the pair \((\mathrm{X} \cup_{f} \mathrm{B}, \alpha_{f}(\mathrm{X}))\) .
\item
  Suppose that the map f is injective and strict. If the pair \((\mathrm{X}, f(\mathrm{A}))\) has the homotopy extension property, so does the pair \((\mathrm{X} \cup_{f} \mathrm{B}, \beta_{f}(\mathrm{B}))\) .
\item
  Suppose the pair (B, A) has the homotopy extension property. Let Z be a topological space, let \(u: X \cup_{f} B \to Z\) be a continuous map and let \(\sigma: \alpha_{f}(X) \times I \to Z\) be a homotopy whose term is the restriction of u to the subspace \(\alpha_{f}(X)\).
\end{enumerate}

Let us set \(v_{1} = u \circ \alpha_{f}\) and denote by \(\tau_{1} \colon \mathbf{X} \times \mathbf{I} \to \mathbf{Z}\) the map given by \((x,t) \mapsto \sigma(\alpha_{f}(x),t)\). Set \(v_{2} = u \circ \beta_{f}\). Since \(\beta_{f} \mid \mathrm{A} = \alpha_{f} \circ f\), the map from \(\mathrm{A} \times \mathbf{I}\) into \(\mathrm{Z}\) which maps \((a,t)\) onto \(\sigma(\alpha_{f}(f(a)),t)\) for \(a \in \mathrm{A}\) and \(t \in \mathbf{I}\) is a homotopy whose term is equal to the map \(v_{2} \mid \mathrm{A}\). Since the pair \((\mathrm{B},\mathrm{A})\) has the homotopy extension property, there exists a homotopy \(\tau_{2} \colon \mathbf{B} \times \mathbf{I} \to \mathbf{Z}\) whose term is the map \(v_{2}\) and such that \(\tau_{2}(a,t) = \sigma(\alpha_{f}(f(a)),t)\) for \((a,t) \in \mathrm{A} \times \mathbf{I}\).

For \(a \in \mathbf{A}\) and \(t \in \mathbf{I}\), we have \(\tau_2(a, t) = \tau_1(f(a), t)\). By Proposition 8, there exists a unique continuous map \(\sigma: (\mathrm{X} \cup_f \mathrm{B}) \times \mathbf{I} \to \mathbf{Z}\) such that \(\sigma(\alpha_f(x), t) = \sigma_1(x, t)\) and \(\sigma(\beta_f(b), t) = \sigma_2(b, t)\), for \(x \in \mathrm{X}\), \(b \in \mathrm{B}\), and \(t \in \mathbf{I}\). This is a homotopy with term \(u\) which extends \(\tau\), whence assertion a).

\begin{enumerate}
\def\labelenumi{\alph{enumi})}
\setcounter{enumi}{1}
\tightlist
\item
  In view of Example 2 (III, p. 249), assertion b) follows from assertion a).
\end{enumerate}

\subsection{Space obtained by contraction of a subspace}

Let X be a topological space and let A be a subset of X. Consider the equivalence relation R in X which is the finest for which all elements of A are equivalent: two elements of X are equivalent under this relation if they are equal or if they both belong to A.

DEFINITION 10. --- The quotient space of X by R is denoted X/A and is called the space obtained from X by contracting the subset A.

Denote by \(\rho\colon X\to X/A\) the canonical surjection. Let Y be a topological space and let \(f\colon X\to Y\) be a continuous map. For there to exist a continuous map \(g\colon X/A\to Y\) such that \(g\circ\rho=f\) , it is necessary and sufficient that \(f\) be constant on A.

If the set A is closed (resp. open) in X, the map \(\rho\) is closed (resp. open) and induces a homeomorphism of \(X\setminus A\) onto its image. If A is a closed and quasi-compact subset of X, the map \(\rho\) is proper (TG, I, p. 75, th. 1).

If the set A is empty, \(\rho\) is a homeomorphism. If it is not empty, A is a point of X/A which is called the base point of X/A and is denoted \(s_{\mathrm{X / A}}\) . The space X/A is then identified with the space obtained by attaching to the space \(\{s_{\mathrm{X / A}}\}\) the space X by means of the unique map from A to \(\{s_{\mathrm{X / A}}\}\) .

PROPOSITION 10. --- Let X be a topological space and let A be a subset of X. Suppose that there exists a homotopy \(\sigma: \mathrm{X} \times \mathbf{I} \to \mathrm{X}\) with source \(\mathrm{Id}_{\mathrm{X}}\), constant on \(\mathrm{A} \times \{1\}\) and such that \(\sigma(\mathrm{A} \times \mathbf{I}) \subset \mathrm{A}\). Then the canonical map \(\rho\) from X to X/A is a homotopy equivalence.

Denote by \(f\) the term of \(\sigma\). This is a continuous map from X to X, homotopic to \(\mathrm{Id}_{\mathrm{X}}\). It is constant on A; hence there exists a unique continuous map \(g \colon \mathrm{X} / \mathrm{A} \to \mathrm{X}\) such that \(g \circ \rho = f\). On the other hand, since \(\sigma(\mathrm{A} \times \mathbf{I}) \subset \mathrm{A}\), there exists a map \(\sigma'\) from \(\mathrm{X} / \mathrm{A} \times \mathbf{I}\) into \(\mathrm{X} / \mathrm{A}\) such that \(\sigma'(\rho(x), t) = \rho(\sigma(x, t))\) for all \(x \in \mathrm{X}\) and every \(t \in \mathbf{I}\). By I, p. 19, prop. 10, the map \(\sigma'\) is continuous. It is a homotopy connecting the identity map of \(\mathrm{X} / \mathrm{A}\) to the map \(\rho \circ g\). Consequently, the maps \(g\) and \(\rho\) are homotopy equivalences inverse to each other up to homotopy.

Remarks. --- Let X be a topological space and let A be a subset of X.

\begin{enumerate}
\def\labelenumi{\arabic{enumi})}
\item
  If the canonical map from X onto X/A is a homotopy equivalence, then there exists a homotopy \(\sigma\colon\mathrm{X}\times\mathbf{I}\to\mathrm{X}\) with source \(\mathrm{Id}_{\mathrm{X}}\) which is constant on \(\mathrm{A}\times\{1\}\). But it may happen that none exists satisfying in addition the condition \(\sigma(\mathrm{A}\times\mathbf{I})\subset\mathrm{A}\) (III, p. 322, exerc. 4).
\item
  There may exist a homotopy connecting the identity map of X to a map from X into X that is constant on A without the spaces X and X/A being homotopy equivalent (III, p. 325, exerc. 14).
\item
  Suppose that A is contractible. If the pair (X, A) has the homotopy extension property, then the canonical map from X to X/A is a homotopy equivalence. Indeed, let \(\sigma\colon A \times I \to A\) be a homotopy whose initial map is the identity map of A and whose terminal map is a constant map. There then exists a homotopy \(\sigma'\) with origin \(\mathrm{Id}_{X}\) which extends \(\sigma\). The assertion thus follows from Proposition 10.
\item
  Let X and Y be topological spaces, let A be a nonempty subspace of X, and let B be a subspace of Y. Let \(f: X \to Y\) be a continuous map such that \(f(A) \subset B\). Denote by \(\varphi: X/A \to Y/B\) the continuous map induced by f on the quotients. If f defines a homotopy equivalence of the pair (X, A) onto the pair (Y, B), then the map \(\varphi\) is a homotopy equivalence of the pair (X/A, \(s_{X/A}\)) onto the pair (Y/B, \(s_{Y/B}\)). Let P be a topological space reduced to a single point, let \(\alpha\) and \(\beta\) be the constant maps from A and B into P. Observe that the map \(\varphi\) is identified with the map from \(P \cup_{\alpha} X\) into \(P \cup_{\beta} Y\) induced by f and the identity map of P. The assertion then follows from Corollary 1 of III, p. 249.
\end{enumerate}

\subsection{Cone of a map}

Let X and Y be topological spaces, and let f be a continuous map from X to Y. Let \(\operatorname{Cyl}(f)\) be the cylinder of the map f. Let us denote \(\alpha_{f}\) the canonical map from \(X \times I\) into \(\operatorname{Cyl}(f)\) and \(f_{0}\) the map \(x \mapsto \alpha_{f}(x,0)\) from X to \(\operatorname{Cyl}(f)\); these maps induce homeomorphisms from \(X \times I\) and X respectively onto their images in \(\operatorname{Cyl}(f)\).

DEFINITION 11. --- The cone of the map f is denoted by \(\operatorname{Côn}(f)\) and is the topological space obtained from \(\operatorname{Cyl}(f)\) by contracting \(f_{0}(\mathrm{X})\).

Denote by \(\beta_{f}^{\prime}\colon\mathrm{Y}\to\mathrm{C}\hat{\mathrm{o}}\mathrm{n}(f)\) the composition of the canonical map \(\beta_{f}\colon\mathrm{Y}\to\mathrm{C}\mathrm{y}\mathrm{l}(f)\) and the canonical surjection from \(\mathrm{C}\mathrm{y}\mathrm{l}(f)\) onto \(\mathrm{C}\hat{\mathrm{o}}\mathrm{n}(f)\). The map \(\beta_{f}^{\prime}\) is continuous and defines a homeomorphism from Y onto a closed subset of \(\mathrm{C}\hat{\mathrm{o}}\mathrm{n}(f)\), called the base of the cone, which we shall thereby identify with Y.

Denote by \(\alpha_{f}^{\prime}\colon\mathrm{X}\times\mathbf{I}\to\mathrm{C}\hat{\mathrm{o}}\mathrm{n}(f)\) the composition of the map \(\alpha_{f}\) and the canonical surjection from \(\mathrm{Cyl}(f)\) onto \(\mathrm{C}\hat{\mathrm{o}}\mathrm{n}(f)\) . The map \(\alpha_{f}^{\prime}\) is a homotopy whose initial map is a constant map and whose terminal map is the map \(\beta_{f}^{\prime}\circ f\) .

If X is empty, the map \(\beta_{f}^{\prime}\) is a homeomorphism from Y onto \(\operatorname{Côn}(f)\) .

Suppose that X is not empty. Let s denote the base point of the space \(\mathrm{Cyl}(f)/f_{0}(\mathrm{X})\); it is called the vertex of the cone \(\mathrm{Côn}(f)\). Since \(f_{0}(\mathrm{X})\) is closed in \(\mathrm{Cyl}(f)\), the canonical map \(\pi\colon\mathrm{Cyl}(f)\to\mathrm{Côn}(f)\) induces a homeomorphism of \(\mathrm{Cyl}(f)\setminus f_{0}(\mathrm{X})\) onto \(\mathrm{Côn}(f)\setminus\{s\}\) (III, p. 252). The vertex of the cone \(\mathrm{Côn}(f)\) does not belong to its base; the canonical injection of Y into \(\mathrm{Côn}(f)\setminus\{s\}\) is a homotopy equivalence (III, p. 239, prop. 6).

Let \(\sigma_f\colon \mathrm{Cyl}(f)\times \mathbf{I}\to \mathrm{Cyl}(f)\) be the canonical contraction of the cylinder of \(f\) onto its base. For \(c\in \mathrm{Cyl}(f)\setminus f_0(\mathrm{X})\) and \(t\in \mathbf{I}\), we have \(\sigma_f(c,t)\notin f_0(\mathrm{X})\). Consequently, there exists a unique map \(\sigma_f'\colon (\mathrm{Côn}(f)\setminus\{s\})\times \mathbf{I}\to \mathrm{Côn}(f)\setminus\{s\}\) such that \(\sigma_f'(\pi(c),t)=\pi(\sigma_f(c,t))\) for \(c\in \mathrm{Cyl}(f)\setminus f_0(\mathrm{X})\) and \(t\in \mathbf{I}\). It is continuous and is a strong contraction of \(\mathrm{Côn}(f)\setminus\{s\}\) onto Y. It is called the canonical contraction of \(\mathrm{Côn}(f)\setminus\{s\}\) onto its base. Its terminal map is a retraction of \(\mathrm{Côn}(f)\setminus\{s\}\) onto Y, called the canonical retraction of the cone with its vertex removed onto its base.

PROPOSITION 11 (Universal property of cones)

Let Z be a topological space, let \(\beta: \mathrm{Y} \to \mathrm{Z}\) be a continuous map and let \(\alpha: \mathrm{X} \times \mathbf{I} \to \mathrm{Z}\) be a homotopy whose initial map is a constant map and whose terminal map is equal to \(\beta \circ f\). There exists a unique continuous map \(\varphi\) from \(\operatorname{Côn}(f)\) into Z such that \(\alpha = \varphi \circ \alpha_f'\) and \(\beta = \varphi \circ \beta_f'\).

By the universal property of cylinders (III, p. 238, prop. 4), there exists a unique continuous map \(h \colon \operatorname{Cyl}(f) \to \mathbf{Z}\) such that \(\alpha = h \circ \alpha_f\) and \(\beta = h \circ \beta_f\) . Since the initial map of \(\alpha\) is a constant map, the restriction of \(h\) to the subspace \(\alpha_f(\mathbf{X} \times \{0\})\) is constant. There therefore exists a unique continuous map \(\varphi \colon \operatorname{Côn}(f) \to \mathbf{Z}\) such that \(h = \varphi \circ \pi\) , where \(\pi\) denotes the canonical surjection of \(\operatorname{Cyl}(f)\) onto \(\operatorname{Côn}(f)\) . The map \(\varphi\) satisfies \(\alpha = \varphi \circ \alpha_f'\) and \(\beta = \varphi \circ \beta_f'\) and is the only one having these properties, since the images of \(\alpha_f'\) and \(\beta_f'\) cover \(\operatorname{Côn}(f)\) .

Example 1. --- Let X be a topological space. We denote by C(X), and call the cone of the space X, the cone of the map IdX; it is the space obtained from X × I by contracting X × \{0\}. Let \(\pi\) be the canonical map from X × I to C(X); the image C'(X) of X × {[}0, 1{[} under \(\pi\) is called the open cone of the space X.

Suppose that X is not empty. Then the cone C(X) is not empty and its base point is still called the vertex of the cone.

The map \(((x,t),u)\mapsto(x,t(1-u))\) from \((\mathrm{X}\times\mathbf{I})\times\mathbf{I}\) into \(X\times I\) is a homotopy connecting the identity map of \(X\times I\) to the map \((x,t)\mapsto(x,0)\). From this we deduce, by passing to quotient sets, a map \(\sigma\colon\mathrm{C}(\mathrm{X})\times\mathbf{I}\to\mathrm{C}(\mathrm{X})\) such that \(\sigma(\pi(x,t),u)=\pi(x,t(1-u))\) for \(x\in X\), \(t,u\in I\). The map \(\sigma\) is continuous by I, p. 19, prop. 10. The map \(\sigma\) is then a pointed homotopy at s connecting the identity map of \(\mathrm{C}(\mathrm{X})\) to the constant map with image \(\{s\}\). Consequently, the cone \(\mathrm{C}(\mathrm{X})\) is contractible at its vertex s. Since \(\sigma(\mathrm{C}'(\mathrm{X})\times\mathbf{I})\) is contained in \(\mathrm{C}'(\mathrm{X})\), the map \(\sigma\) defines, by passing to subspaces, a pointed homotopy at s connecting the identity map of \(\mathrm{C}'(\mathrm{X})\) to the constant map with image \(\{s\}\); the open cone \(\mathrm{C}'(\mathrm{X})\) is therefore contractible at s.

Remarks. --- 1) Let X be a topological space and let A be a subspace of X; denote by i the canonical injection of A into X. The canonical retraction \(r_{i}\) of the cylinder \(\mathrm{Cyl}(i)\) onto its base maps the subspace \(\alpha_{i}(\mathrm{A} \times \{0\})\) into A. Let \(\rho: \mathrm{C}\hat{\mathrm{o}}\mathrm{n}(i) \to \mathrm{X}/\mathrm{A}\) be the continuous map induced therefrom by passing to quotients. Suppose that the pair (X, A) has the homotopy extension property. By III, p. 243, Theorem 1, the map \(r_{i}\) defines a homotopy equivalence of the pair \((\mathrm{Cyl}(i), \alpha_{i}(\mathrm{A} \times \{0\}))\) onto the pair (X, A). By Remark 4 of III, p. 253, the map \(\rho\) is then a homotopy equivalence of the pair \((\mathrm{C}\hat{\mathrm{o}}\mathrm{n}(i), s)\) onto the pair \((\mathrm{X}, s_{\mathrm{X}/\mathrm{A}})\). It is in particular a homotopy equivalence.

\begin{enumerate}
\def\labelenumi{\arabic{enumi})}
\setcounter{enumi}{1}
\tightlist
\item
  Let X and Y be topological spaces, and let \(f\colon X \to Y\) be a continuous map. The canonical map \(\alpha_{f}^{\prime}\colon X \times I \to \operatorname{Côn}(f)\) is constant on \(X \times \{0\}\) and therefore defines a continuous map \(\gamma_{f}\colon C(X) \to \operatorname{Côn}(f)\). The restriction to \(C'(X)\) of the map \(\gamma_{f}\) is injective and strict, and induces by passing to subspaces a homeomorphism of the open cone \(C'(X)\) onto the complement of Y in \(\operatorname{Côn}(f)\).
\end{enumerate}

Identify the base of the cone \(C(X)\) with the space X, and denote by \(u\) the continuous map from \(Y \cup_f C(X)\) into \(\operatorname{Côn}(f)\) induced by the maps \(\beta_f':Y\to\operatorname{Côn}(f)\) and \(\gamma_f:C(X)\to\operatorname{Côn}(f)\). Conversely, let \(v:\operatorname{Côn}(f)\to Y\cup_f C(X)\) be the unique continuous map such that \(v\circ\alpha_f'\) is the canonical map from \(X\times I\) onto \(C(X)\), and \(v\circ\beta_f'\) is the canonical map from \(Y\) to \(Y\cup_f C(X)\). The maps \(u\) and \(v\) are mutually inverse homeomorphisms.

\section{HOMOTOPY AND PATHS}

\subsection{Paths}

DEFINITION 1. --- Let X be a topological space. A path in X is any continuous map c from I to X. The point \(c(0)\) is called the origin, the point \(c(1)\) the endpoint of the path c. A loop in X is a path in X whose origin is equal to the endpoint.

Let \(x\) and \(y\) be points of X. We say that a path \(c\) in X connects \(x\) to \(y\) if \(x\) is its origin and \(y\) its endpoint.

DEFINITION 2. --- Let X be a topological space. Paths c and d in X are said to be juxtaposable if one has \(c(1) = d(0)\). We then call the juxtaposed path of c and d the path \(c*d\) defined by the formula:

\[
(c * d) (t) = \left\{ \begin{array}{l l} c (2 t) & \text { for } 0 \leqslant t \leqslant 1 / 2, \\ d (2 t - 1) & \text { for } 1 / 2 \leqslant t \leqslant 1. \end{array} \right.\tag{1}
\]

Its origin is the origin of c, and its endpoint is that of d.

Let c be a path in X; one calls the path opposite to c, and denotes by \(\overline{c}\) the path defined by \(\overline{c}(t) = c(1 - t)\) for \(t \in I\).

If c and d are two juxtaposable paths in X, \(\overline{d}\) and \(\overline{c}\) are, and one has \(\overline{c*d} = \overline{d} * \overline{c}\) . For every path c in X, one has \(\overline{\overline{c}} = c\) .

Remarks. --- 1) Let X be a topological space and let P be a topological space reduced to a point. Identify P × I with I by the projection pr₂. The set C(P × I; X) of homotopies between maps (necessarily continuous) from P to X is then identified with the set C(I; X) of paths in X. To a homotopy connecting the constant map with image x to the constant map with image y there corresponds a path with origin x and endpoint y. The preceding identifications are compatible with the notions of juxtaposition and passage to the opposite (cf. III, p. 230).

\begin{enumerate}
\def\labelenumi{\arabic{enumi})}
\setcounter{enumi}{1}
\tightlist
\item
  Let X and Y be topological spaces. The canonical map
\end{enumerate}

\[
\mathcal {C} (\mathrm{X} \times \mathbf {I}; \mathrm{Y}) \rightarrow \mathcal {C} (\mathbf {I}; \mathcal {C} _ {c} (\mathrm{X}; \mathrm{Y}))
\]

associates to every homotopy \(\sigma\colon X \times I \to Y\) connecting two maps \(f\) and \(g\) from X to Y a path of origin \(f\), with endpoint \(g\), in the space \(\mathcal{C}_c(X;Y)\) . If \(X\) is a locally compact space, this canonical map is bijective (TG, X, p. 28, theorem 3).

Let X be a topological space. One calls the path space of X, and denotes by \(\Lambda(\mathrm{X})\) , the topological space \(\mathcal{C}_{c}(\mathbf{I};\mathrm{X})\) whose elements are the paths in X and whose topology is that of compact convergence (cf. TG, X, p. 27). If x is a point of X, one denotes by \(\Lambda_{x}(\mathrm{X})\) the subspace of \(\Lambda(\mathrm{X})\) formed by the paths with origin x. If y is a second point of X, one denotes by \(\Lambda_{x,y}(\mathrm{X})\) the subspace of \(\Lambda(\mathrm{X})\) formed by the paths with origin x and endpoint y.

PROPOSITION 1. --- Let X and Z be topological spaces. For a map \(\varphi\) from Z into the path space \(\mathcal{C}_c(\mathbf{I};\mathrm{X})\) to be continuous, it is necessary and sufficient that the map \((z,t)\mapsto \varphi (z)(t)\) be a continuous map from \(\mathbf{Z}\times \mathbf{I}\) into X. In particular, the map \((c,t)\mapsto c(t)\) of \(\mathcal{C}_c(\mathbf{I};\mathrm{X})\times \mathbf{I}\) into X is continuous.

The topological space I being locally compact, the proposition follows from TG, X, p. 28, th. 3.

One denotes by o and calls the origin map the map \(c \mapsto c(0)\) of \(\Lambda(\mathrm{X})\) into X; one denotes by e and calls the endpoint map the map \(c \mapsto c(1)\) of \(\Lambda(\mathrm{X})\) into X. The maps o and e are continuous (TG, X, p. 27, remark 1). The pairs of juxtaposable paths in X are the elements of the fibered product of the X-spaces \((\Lambda(\mathrm{X}), e)\) and \((\Lambda(\mathrm{X}), o)\) .

PROPOSITION 2. --- Let X be a topological space. The map which to a path c associates the opposite path \(\overline{c}\) is a homeomorphism of \(\Lambda(X)\) on itself. The juxtaposition of paths \((c,d)\mapsto c*d\) is a homeomorphism of the space \((\Lambda(X),e)\times_{X}(\Lambda(X),o)\) onto \(\Lambda(X)\) .

The map \(t \mapsto 1 - t\) is a homeomorphism of the space I onto itself; the first assertion follows from this.

Let C denote the fiber product \((\Lambda(\mathrm{X}),e)\times_{\mathrm{X}}(\Lambda(\mathrm{X}),o)\) . Denote by \(\gamma:C\to\Lambda(X)\) the map \((c,d)\mapsto c*d\) and \(\delta:C\times I\to X\) the map \(((c,d),t)\mapsto(c*d)(t)\) . The restrictions of the map \(\delta\) to \(C\times[0,\frac{1}{2}]\) and \(C\times[\frac{1}{2},1]\) are continuous by virtue of formula (1) and prop. 1. The map \(\delta\) is therefore continuous (TG, I, p. 19, prop. 4), as is the map \(\gamma\) (prop. 1). We have \(c(t)=(c*d)(\frac{t}{2})\) and \(d(t)=(c*d)(\frac{1+t}{2})\) , therefore \(\gamma\) is injective.

Let \(g\) be a path in \(\mathbf{X}\); for \(t \in \mathbf{I}\), set

\[
c _ {g} (t) = g \bigl (\frac {t}{2} \bigr) \quad \text{and} \quad d _ {g} (t) = g \bigl (\frac {1 + t}{2} \bigr).
\]

The maps \(c_{g}\) and \(d_{g}\) thus defined are juxtaposable paths in X and we have \(c_{g} * d_{g} = g\) . Moreover, the maps \(g \mapsto c_{g}\) and \(g \mapsto d_{g}\) are continuous maps of the space \(\Lambda(\mathrm{X})\) into itself (prop. 1). It follows that the map \(\gamma\) is a homeomorphism.

COROLLARY. --- Let X be a topological space and let x, y, z be points of X. The map \(c \mapsto \overline{c}\) from \(\Lambda_{x,y}(X)\) to \(\Lambda_{y,x}(X)\) and the map \((c,d) \mapsto c * d\) from \(\Lambda_{x,y}(X) \times \Lambda_{y,z}(X)\) to \(\Lambda_{x,z}(X)\) are continuous.

\subsection{Path-connected spaces}

DEFINITION 3. --- A topological space X is said to be path-connected if for every pair \((x,y)\) of points of X, there exists a path with origin x and endpoint y.

Example 1. --- Every interval of R, every convex subset of a numerical space or, more generally, of a topological vector space over R is path-connected.

PROPOSITION 3. --- The image under a continuous map of a path-connected space is path-connected.

Let X be a path-connected space, \(f: X \to Y\) a continuous map. Let x and y be two points of \(f(X)\); let \(x'\) and \(y'\) be points of X such that \(f(x') = x\) and \(f(y') = y\). There exists a path c in X whose origin is \(x'\) and endpoint \(y'\). Then, \(f \circ c\) is a path in \(f(X)\) with origin x and endpoint y.

PROPOSITION 4. --- A path-connected topological space is connected.

Since every interval of R is connected, the image of a path is connected (TG, I, p. 82, prop. 4). The empty space is connected. Since a nonempty path-connected space is the union of the images of paths issuing from one of its points, it is connected (TG, I, p. 81, prop. 2).

In view of the identification (III, p. 256, remark 1) of paths in X with homotopies between maps of a topological space P reduced to a point into X, proposition 1 of III, p. 230 implies:

PROPOSITION 5. --- In a topological space, the relation “there exists a path with origin x and endpoint y” is an equivalence relation.

The path-connected components of a topological space X are called the equivalence classes for the above relation. A path-connected component is a path-connected space. Let x be a point of X. The path-connected component of x is the union of the path-connected subspaces of X containing x. It is also the union of the images of paths with origin x in X.

Examples. --- 2) The union of a family of path-connected sets, whose intersection is nonempty, is path-connected (cf. TG, I, p. 81, prop. 2).

\begin{enumerate}
\def\labelenumi{\arabic{enumi})}
\setcounter{enumi}{2}
\tightlist
\item
  A subset of a numerical space (or, more generally, of a topological vector space over R) which is star-shaped at one of its points is path-connected.
\end{enumerate}

Let X be a topological space. We denote by \(\pi_{0}(X)\) the set of path-connected components of X. Let P be a topological space reduced to a point. The map from X to {[}P; X{]} which, to every point x of X, associates the homotopy class \(\varphi_{x}\) of the map \(f_{x}: P \to X\) with image x, defines by passage to the quotient a bijection, called canonical, from \(\pi_{0}(X)\) onto {[}P; X{]}. We have \(\pi_{0}(\varnothing) = \varnothing\) . For a nonempty topological space to be path-connected, it is necessary and sufficient that \(\pi_{0}(X)\) be a set with one element.

Let X and Y be topological spaces and \(f\colon\mathrm{X}\to\mathrm{Y}\) a continuous map. The image \(f(\mathrm{C})\) of every path-connected component C of X is path-connected (III, p. 258, prop. 3), hence contained in a path-connected component of Y. We denote by \(\pi_{0}(f)\colon\pi_{0}(\mathrm{X})\to\pi_{0}(\mathrm{Y})\) the map which, to a path-connected component C of X, associates the unique path-connected component C' of Y such that \(f(\mathrm{C})\subset\mathrm{C}'\). If one identifies \(\pi_0(\mathrm{X})\) and \(\pi_0(\mathrm{Y})\) with {[}P; X{]} and {[}P; Y{]} respectively, \(\pi_0(f)\) is identified with the map \(\chi \mapsto [f] \circ \chi\) from {[}P; X{]} to {[}P; Y{]}. In particular, if \(f\) and \(g\) are homotopic maps from X to Y, we have \(\pi_0(f) = \pi_0(g)\) (III, p. 230, prop. 2).

Let Z be a topological space and let \(g: Y \to Z\) be a continuous map; we have \(\pi_{0}(g \circ f) = \pi_{0}(g) \circ \pi_{0}(f)\) . In particular, if Z = X and if f and g are homotopy inverses of each other, we have \(\pi_{0}(g) \circ \pi_{0}(f) = \pi_{0}(\mathrm{Id}_{\mathrm{X}}) = \mathrm{Id}_{\pi_{0}(\mathrm{X})}\) and \(\pi_{0}(f) \circ \pi_{0}(g) = \pi_{0}(\mathrm{Id}_{\mathrm{Y}}) = \mathrm{Id}_{\pi_{0}(\mathrm{Y})}\) , which proves that \(\pi_{0}(f)\) and \(\pi_{0}(g)\) are inverse bijections of each other. A space homotopy equivalent to a path-connected space is therefore itself path-connected. In particular, a space homotopy equivalent to a point is path-connected.

PROPOSITION 6. --- Let \((\mathrm{Y}_{j})_{j\in\mathrm{J}}\) be a family of topological spaces. The map

\[
(\pi_ {0} (\mathrm{pr} _ {j})) \colon \pi_ {0} (\prod_ {j \in \mathrm{J}} \mathrm{Y} _ {j}) \to \prod_ {j \in \mathrm{J}} \pi_ {0} (\mathrm{Y} _ {j})
\]

is bijective. In particular, the product space of a family of path-connected spaces is path-connected.

This follows from Proposition 3 of III, p. 231, where one takes for X a space reduced to a single point.

\subsection{Locally path-connected spaces}

DEFINITION 4. --- A topological space is said to be locally path-connected if each of its points has a fundamental system of path-connected neighborhoods.

PROPOSITION 7. --- A locally path-connected topological space is locally connected. Its connected components coincide with its path-connected components. In particular, if a locally path-connected topological space is connected, it is path-connected.

The first assertion follows from Prop. 4. Let us prove the second assertion. Let X be a locally path-connected topological space and let C be a path-connected component of X. Every point of C has a path-connected neighborhood, hence contained in C; consequently, C is an open subset of X. Since the path-connected components form a partition of X, such a component is also closed. As it is connected (III, p. 258, prop. 4), it is a connected component (TG, I, p. 83). The third assertion follows from the second.

Note thus that there is no ambiguity in saying that a topological space is connected and locally path-connected.

COROLLARY 1. --- Every connected open subset of a locally path-connected topological space is path-connected.

COROLLARY 2. --- Let B be a locally path-connected topological space and let E be an étale B-space. The space E is locally path-connected. If it is connected, it is path-connected.

The first assertion follows immediately from def. 4 and the definition of an étale morphism (I, p. 28, def. 2). The second is deduced from it by prop. 7.

For a topological space X to be locally path-connected, it is necessary and sufficient that every path-connected component of an open subset of X be an open subset of X (cf. I, p. 85, proof of Prop. 11). If the space X is locally path-connected, then every point of X has a fundamental system of open path-connected neighborhoods.

PROPOSITION 8. --- Every quotient space of a locally path-connected space is locally path-connected.

Let X be a locally path-connected space, let R be an equivalence relation on X, and let \(\varphi\colon\mathrm{X}\to\mathrm{X}/\mathrm{R}\) be the canonical map. It suffices to show that the path-connected components of an open subset of X/R are open. Let A thus be an open subset of X/R and let C be a path-connected component of A. Let \(x\in\bar{\varphi}^{1}(\mathrm{C})\), and let K be the path-connected component of x in \(\bar{\varphi}^{1}(\mathrm{A})\). The set \(\varphi(\mathrm{K})\) is path-connected (III, p. 258, prop. 3), contained in A and contains \(\varphi(x)\); one therefore has \(\varphi(\mathrm{K})\subset\mathrm{C}\), whence \(\mathrm{K}\subset\bar{\varphi}^{1}(\mathrm{C})\). This proves that \(\bar{\varphi}^{1}(\mathrm{C})\) is a union of path-connected components of \(\bar{\varphi}^{1}(\mathrm{A})\). Since X is locally path-connected and \(\bar{\varphi}^{1}(\mathrm{A})\) is open in X, \(\bar{\varphi}^{1}(\mathrm{C})\) is open in X. Consequently, C is open in X/R. This proves the proposition.

PROPOSITION 9. --- The product of a family of locally path-connected spaces, all but finitely many of which are connected, is locally path-connected.

Let \((\mathrm{X}_{j})_{j\in\mathrm{J}}\) be a family of locally path-connected spaces that are connected, except for finitely many of them. Let \(x=(x_{j})_{j\in\mathrm{J}}\) be a point of the product space \(X=\prod_{j\in J}X_{j}\). By hypothesis, a fundamental system of neighborhoods of x in X consists of the sets of the form \(V=\prod_{j\in J}V_{j}\) where, for every \(j\in J\), \(V_{j}\) is a path-connected neighborhood of \(x_{j}\) in \(X_{j}\) and where \(V_{j}=X_{j}\) except for a finite set of indices \(j\in J\) (TG, I, p. 24). These sets being path-connected (III, p. 260, prop. 6), X is locally path-connected.

COROLLARY. --- Every open subset of a numerical space is locally path-connected.

Every point of R has a base of neighborhoods consisting of intervals; consequently, R is locally path-connected. By Proposition 9, the numerical space \(R^{n}\) is locally path-connected, for every integer \(n \geqslant 1\). Furthermore, an open subset of a locally path-connected topological space is still locally path-connected, whence the corollary.

Example. --- Let G be the subgroup of \(\mathbf{GL}(n,\mathbf{R})\) consisting of square matrices of order n whose determinant is strictly positive. The group G is connected and locally path-connected.

By A, III, p. 104, prop. 17, the group \(\mathbf{SL}(n,\mathbf{R})\) is generated by the elements \(B_{ij}(\lambda)\) (for \(1\leqslant i,j\leqslant n\) such that \(i\neq j\) and \(\lambda\in\mathbf{R}\)). The maps \(\lambda\mapsto B_{ij}(\lambda)\) from \(\mathbf{R}\) into \(\mathbf{GL}(n,\mathbf{R})\) are continuous, their images are connected subsets of \(\mathbf{SL}(n,\mathbf{R})\); since they all contain the identity matrix \(I_n\) their union is connected (TG, I, p. 81, prop. 2). Consequently, the group \(\mathbf{SL}(n,\mathbf{R})\) is connected (TG, III, p. 8, prop. 7). Let A be the subgroup of G formed by the matrices of the form diag(1, \(\ldots\),1, \(\lambda\)), with \(\lambda\in\mathbf{R}_{+}^{*}\); it is connected and one has \(\mathbf{GL}(n,\mathbf{R})=\mathbf{A}\cdot\mathbf{SL}(n,\mathbf{R})\). It follows that the group G is connected. Since it is the inverse image of \(\mathbf{R}_{+}^{*}\) under the determinant map from \(\mathbf{M}_{n}(\mathbf{R})\) into \(\mathbf{R}\), it is an open subset of \(\mathbf{M}_{n}(\mathbf{R})\); it is therefore locally path-connected (corollary above), as well as path-connected (III, p. 261, corollary 1).

PROPOSITION 10. --- Let X be a topological space. The map \((o, e)\) from \(\Lambda(\mathrm{X}) = \mathcal{C}_{c}(\mathbf{I}; \mathrm{X})\) into \(X \times X\), which associates to a path c in X the pair \((c(0), c(1))\), is continuous. If the space X is locally path-connected, this map is open.

We already know that the origin map o and the endpoint map e of \(\Lambda(\mathrm{X})\) into X are continuous (III, p. 257), whence the first assertion.

Lemma. --- Let \(c\colon \mathbf{I}\to \mathrm{X}\) be a path and let W be a neighborhood of \(c\) in the space \(\mathcal{C}_c(\mathbf{I};\mathrm{X})\). There exists a real number \(\varepsilon \in ]0,1 / 2]\), a neighborhood \(\mathrm{V}_0\) of \(c(0)\) and a neighbourhood \(\mathrm{V}_1\) of \(c(1)\) in X having the following properties: we have \(c([0,\varepsilon])\subset \mathrm{V}_0\), \(c([1 - \varepsilon,1])\subset \mathrm{V}_1\), and every path \(c^{\prime}\colon \mathbf{I}\rightarrow \mathrm{X}\) such that

\[
\left\{ \begin{array}{l l} c ^ {\prime} (t) \in \mathrm{V} _ {0} & \text { for } 0 \leqslant t \leqslant \varepsilon , \\ c ^ {\prime} (t) = c (t) & \text { for } \varepsilon \leqslant t \leqslant 1 - \varepsilon , \\ c ^ {\prime} (t) \in \mathrm{V} _ {1} & \text { for } 1 - \varepsilon \leqslant t \leqslant 1, \end{array} \right.\tag{2}
\]

belongs to W.

By definition of the topology of compact convergence (TG, X, p. 26, def. 1), there exists a finite set J, a family \((\mathrm{U}_{j})_{j\in\mathrm{J}}\) of open sets in X and a family \((\mathrm{K}_{j})_{j\in\mathrm{J}}\) of compact subsets of I such that the set \(W'\) of paths \(c'\) satisfying \(c'(\mathrm{K}_{j})\subset\mathrm{U}_{j}\) for every index j is a neighborhood of c contained in W. Let us then denote by \(A_{0}\) (resp. \(A_{1}\) ) the set of indices j such that \(0\in K_{j}\) (resp. \(1\in K_{j}\) ); set \(V_{0}=\bigcap_{j\in A_{0}}U_{j}\) and \(V_{1}=\bigcap_{j\in A_{1}}U_{j}\) .

Since the map \(c\) is continuous, there exists a real number \(\varepsilon \in [0,1/2]\) such that \(c([0,\varepsilon]) \subset \mathrm{V}_0\) , \(c([1-\varepsilon,1]) \subset \mathrm{V}_1\) , \([0,\varepsilon] \cap \mathrm{K}_j = \emptyset\) for all \(j \notin \mathrm{A}_0\) and \([1-\varepsilon,1] \cap \mathrm{K}_j = \emptyset\) for all \(j \notin \mathrm{A}_1\) Then let \(c'\) be a path satisfying conditions (2). Let us prove that \(c' \in \mathrm{W}'\) . Let \(j \in \mathrm{J}\) and let \(t \in \mathrm{K}_j\) . If \(\varepsilon \leqslant t \leqslant 1-\varepsilon\) , \(c'(t) = c(t)\) belongs to \(\mathrm{U}_j\) . If \(0 \leqslant t \leqslant \varepsilon\) , \(c'(t) \in \mathrm{V}_0\) ; by the choice of \(\varepsilon\) , we have \(j \in \mathrm{A}_0\) , therefore \(c'(t) \in \mathrm{U}_j\) . Similarly, if \(1-\varepsilon \leqslant t \leqslant 1\) , we have \(j \in \mathrm{A}_1\) and \(c'(t) \in \mathrm{V}_1 \subset \mathrm{U}_j\) . Thus, \(c'(\mathrm{K}_j) \subset \mathrm{U}_j\) and \(c'\) belongs to \(\mathrm{W}'\) , hence to W.

Let us now prove the second assertion of prop. 10. Suppose the space X is locally path-connected. Let \(c\) be a path in X and let W be a neighborhood of \(c\) in \(\mathcal{C}_c(\mathbf{I};\mathbf{X})\). Let \(\varepsilon\), \(\mathrm{V}_0\), and \(\mathrm{V}_1\) be as in the lemma. Let \(\mathrm{T}_0\) and \(\mathrm{T}_1\) be path-connected neighborhoods of \(c(0)\) and \(c(1)\) contained in \(\mathrm{V}_0\) and \(\mathrm{V}_1\), respectively. There exists a real number \(\theta\) such that \(0 < \theta < \varepsilon\) and such that \(c([0,\theta]) \subset \mathrm{T}_0\), \(c([1 - \theta,1]) \subset \mathrm{T}_1\). Let \(x_0 \in \mathrm{T}_0\) and \(x_1 \in \mathrm{T}_1\); let \(c_0\) be a path of origin \(x_0\) and of endpoint \(c(\theta)\) in \(\mathrm{T}_0\) and let \(c_1\) be a path of origin \(x_1\) and of endpoint \(c(1 - \theta)\)

in \(\mathrm{T}_1\). Set

\[
c ^ {\prime} (t) = \left\{ \begin{array}{l l} c _ {0} (t / \theta) & \text { for } 0 \leqslant t \leqslant \theta , \\ c (t) & \text { for } \theta \leqslant t \leqslant 1 - \theta , \\ c _ {1} ((1 - t) / \theta) & \text { for } 1 - \theta \leqslant t \leqslant 1. \end{array} \right.
\]

One thus defines a path \(c'\) connecting \(x_{0}\) to \(x_{1}\) and satisfying conditions (2). This proves that the image of W in X × X under the map \((o, e)\) contains the neighborhood \(T_{0} \times T_{1}\) of \((c(0), c(1))\) , whence the proposition.

COROLLARY. --- Let X be a locally path-connected topological space and let x be a point of X. The map \(c \mapsto c(1)\) of \(\Lambda_{x}(\mathrm{X})\) into X is open.

By the proposition, the map \(\varphi\colon\Lambda(\mathrm{X})\to\mathrm{X}\times\mathrm{X}\) defined by \(\varphi(c)=(c(0),c(1))\) is open and one has \(\Lambda_{x}(\mathrm{X})=\overline{\varphi}^{1}(\{x\}\times\mathrm{X})\) . Consequently, the map \(\Lambda_{x}(\mathrm{X})\to\{x\}\times\mathrm{X}\) deduced from \(\varphi\) is open (TG, I, p. 30, prop. 2), as is its composite with the second projection \(\mathrm{pr}_{2}\) .

\subsection{Links between connectedness and path connectedness}

A path-connected space is connected (III, p. 258, prop. 4). There exist connected spaces, even locally connected ones, that are not path-connected (cf. III, p. 331, exerc. 2 and 4). However:

PROPOSITION 11. --- A connected and locally connected topological space, whose topology can be defined by a distance for which it is complete, is path-connected.

Let X be a topological space. Call a train in X any nonempty finite sequence \(\mathrm{T} = (\mathrm{W}_{i})_{1 \leqslant i \leqslant n}\) of connected open subsets of X such that \(W_{i} \cap W_{i+1} \neq \varnothing\) for \(1 \leqslant i \leqslant n-1\). We say that n is the length of the train T and that the \(W_{i}\) are its cars. If X is equipped with a distance compatible with its topology, we call the width of the train T the maximum of the diameters of its cars. We say that the train joins a point a to a point b if a belongs to the first and b to the last car. We call a refinement of T any pair \((\mathrm{T}', f)\) formed by a train \(\mathrm{T}' = (\mathrm{W}_{j}')_{1 \leqslant j \leqslant m}\) and a strictly increasing map \(f: \{0, 1, \ldots, n\} \to \{0, 1, \ldots, m\}\) such that \(f(0) = 0\), \(f(n) = m\) and \(W_{j}' \subset W_{i}\) for \(1 \leqslant i \leqslant n\) and \(f(i-1) < j \leqslant f(i)\).

Lemma 1. --- Let X be a connected and locally connected metric space, let a and b be points of X and let \(\varepsilon\) be a real number \(>0\) . There exists in X a train of width \(\leqslant \varepsilon\) joining a to b.

More precisely, every train T joining a to b has a refinement \((\mathrm{T}^{\prime}, f)\), where \(T^{\prime}\) is a train of width \(\leqslant\varepsilon\) joining a to b.

The relation “there exists a train of width \(\leqslant\varepsilon\) joining x to y” is an equivalence relation between x and y in X. The equivalence class of a point x contains every connected open neighborhood of x of diameter \(\leqslant\varepsilon\) and x possesses such a neighborhood since X is locally connected. The equivalence classes under this relation are therefore open, and consequently also closed. There is at most one of them, since X is connected. This proves the first assertion.

Let \(\mathrm{T}=(\mathrm{W}_{i})_{1\leqslant i\leqslant n}\) be a train in X joining a to b. Set \(x_{0}=a\), \(x_{n}=b\) and choose for \(1\leqslant i\leqslant n-1\) a point \(x_{i}\) in the nonempty set \(W_{i}\cap W_{i+1}\). For \(1\leqslant i\leqslant n\), the open set \(W_{i}\) is connected and locally connected and \(x_{i-1}\), \(x_{i}\) are two of its points; by the preceding paragraph, there exists a train \((\mathrm{W}_{i,k})_{1\leqslant k\leqslant m_{i}}\) in \(W_{i}\) of width \(\leqslant\varepsilon\), joining \(x_{i-1}\) to \(x_{i}\).

Let us set \(m = m_{1} + \cdots + m_{n}\) . For \(1 \leqslant j \leqslant m\) , set \(W_{j}^{\prime} = W_{i,k}\) , where \((i, k)\) is the unique pair of integers such that \(1 \leqslant i \leqslant n\) , \(1 \leqslant k \leqslant m_{i}\) and \(j = m_{1} + \cdots + m_{i-1} + k\) . For \(0 \leqslant i \leqslant n\) , set \(f(i) = m_{1} + \cdots + m_{i}\) . Then \(T^{\prime} = (W_{j}^{\prime})_{1 \leqslant j \leqslant m}\) is a train of width \(\leqslant \varepsilon\) in X joining a to b, and \((T^{\prime}, f)\) is a refinement of T.

Let us now prove Proposition 11. Endow the connected and locally connected space X with a distance d, compatible with its topology, for which it is complete. Let a and b be points of X. Lemma 1 allows one to construct, by induction, sequences \((\mathrm{T}_s)_{s \geqslant 0}\) and \((f_s)_{s \geqslant 1}\) such that, for all \(s \geqslant 0\), \(\mathrm{T}_s = (\mathrm{W}_{s,i})_{1 \leqslant i \leqslant n_s}\) is a train of width \(\leqslant 2^{-s}\) joining a to b and \((\mathrm{T}_{s+1}, f_{s+1})\) is a refinement of \(\mathrm{T}_s\).

We can choose by induction, for \(s \geqslant 0\), a strictly increasing map \(g_{s} \colon \{0,1,\ldots,n_{s}\} \to I\) such that \(g_{s}(0) = 0\) and \(g_{s}(n_{s}) = 1\), so that \(g_{s+1} \circ f_{s} = g_{s}\). Let us define, for \(s \geqslant 0\), a subset \(A_{s}\) of \(I \times X\) by setting

\[
\mathrm{A} _ {s} = \bigcup_ {1 \leqslant i \leqslant n _ {s}} ([ g _ {s} (i - 1), g _ {s} (i) ] \times \mathrm{W} _ {s, i}).
\]

The sequence \((\mathrm{A}_{s})_{s\geqslant0}\) is decreasing: indeed, for every integer \(j\in\{1,\ldots,n_{s+1}\}\) , there exists a unique integer \(i\in\{1,\ldots,n_{s}\}\) such that \(f_{s}(i-1)<j\leqslant f_{s}(i)\) and one has

\[
[ g _ {s + 1} (j - 1), g _ {s + 1} (j) ] \subset [ g _ {s} (i - 1), g _ {s} (i) ], \quad \mathrm{W} _ {s + 1, j} \subset \mathrm{W} _ {s, i}.
\]

Let \(t \in \mathbf{I}\) . For all \(s \geqslant 0\) , denote by \(\mathrm{A}_s(t)\) the set of \(x \in \mathrm{X}\) such that \((t, x) \in \mathrm{A}_s\) . The set \(\mathrm{A}_s(t)\) is either one of the cars, or the union of two consecutive cars of the train \(\mathrm{T}_s\) ; it is therefore a nonempty subset of \(\mathrm{X}\) , of diameter \(\leqslant 2^{1-s}\) . The sequence \((\mathrm{A}_s(t))_{s \geqslant 0}\) is decreasing. The set of its terms is a Cauchy filter base. This converges to a point \(c(t)\) since the metric space \(\mathrm{X}\) is complete.

Since \(a\) belongs to each of the sets \(\mathrm{A}_s(0) = \mathrm{W}_{s,1}\), we have \(c(0) = a\); similarly, we have \(c(1) = b\). Let \(t \in \mathbf{I}\) and let \(s\) be an integer \(\geqslant 0\). The point \(t\) has in \(\mathbf{I}\) a neighborhood V of one of the following forms: \([g_s(0), g_s(1)[\), \(]g_s(i - 1), g_s(i + 1)[\) for an integer \(i\) such that \(1 \leqslant i \leqslant n_s - 1\), or \(]g_s(n_s - 1), g_s(n_s)]\). Depending on the case, the set \(c(\mathrm{V})\) is contained in the closure of the first car, of the union of two consecutive cars, or of the last car of the train \(\mathrm{T}_s\). It is therefore of diameter \(\leqslant 2^{1-s}\) and one has \(d(c(t), c(t')) \leqslant 2^{1-s}\) for \(t' \in \mathrm{V}\). This proves that the map \(c\colon \mathbf{I} \to \mathbf{X}\) is continuous. Thus \(c\) is a path in X joining \(a\) to \(b\), and X is path-connected.

COROLLARY 1. --- A locally connected topological space whose topology can be defined by a distance for which it is complete is locally path-connected.

Let X be such a space and let U be an open and connected subset of X. By Lemma 2 below, there exists a distance on U compatible with its topology for which U is complete. Since U is locally connected, it follows from Proposition 11 that U is path-connected.

Lemma 2. --- Let X be a complete metric space and let U be an open subset of X. There exists a distance on U compatible with its topology for which U is complete.

We shall take up the arguments of the proof of Prop. 2 of TG, IX, p. 57. We may assume U is distinct from X. Let V be the subset of the product \(\mathbf{R} \times \mathbf{X}\) formed by the points \((t,x)\) such that \(td(x,\mathrm{X}-\mathrm{U})=1\) ; the subspace V of \(\mathbf{R} \times \mathbf{X}\) is closed and the map \((t,x) \mapsto x\) from V to U is a homeomorphism (TG, IX, p. 13, Prop. 3). There exists on \(\mathbf{R} \times \mathbf{X}\) a distance \(d'\) compatible with its topology for which \(\mathbf{R} \times \mathbf{X}\) is complete (TG, IX, p. 15, Cor. 2 and TG, II, p. 17, Prop. 10). The space V is complete for the distance induced by \(d'\) (TG, II, p. 16, Prop. 8), whence the lemma.

COROLLARY 2. --- A locally compact, locally connected, and metrizable topological space is locally path-connected. If it is connected, it is path-connected.

It suffices to prove the first assertion (III, p. 260, Prop. 7). Let X be a locally compact and locally connected metric space. The open, connected, and relatively compact subsets of X form a base for the topology of X. Let U be such a set; since U is open in its closure, which is a compact, hence complete, metric space, there exists by Lemma 2 a distance on U compatible with its topology for which U is complete. It follows from Prop. 11 of III, p. 264 that U is path-connected, whence the corollary.

\subsection{Path-continuous maps}

DEFINITION 5. --- Let X and Y be topological spaces. A map \(f: X \to Y\) is said to be path-continuous if, for every path \(c\) in X, the map \(f \circ c: I \to Y\) is continuous.

Remark. --- Suppose that the space X is path-connected. Let x be a point of X. For f to be path-continuous, it suffices that, for every path c with origin x, the map \(f \circ c\) be continuous. Indeed, let d be any path in X and let c be a path in X with origin x and endpoint \(d(0)\) . If the map \(f \circ (c * d)\) is continuous, then so is the map \(f \circ d\) because we have \(f \circ d(t) = f \circ (c * d)((t + 1)/2)\) .

A continuous map is path-continuous. The converse is not always true; Propositions 12 and 13 below provide criteria for asserting that a path-continuous map is continuous.

PROPOSITION 12. --- Let X and Y be topological spaces and let \(f: \mathrm{X} \to \mathrm{Y}\) be a map. Suppose that the space X is locally path-connected and that every point of X has a countable fundamental system of neighborhoods. If the map \(f\) is path-continuous, it is continuous.

By (TG, IX, p. 18), it suffices to prove that, for every point \(x\) of X and every sequence \((x_n)_{n \geqslant 1}\) of points of X tending to \(x\) , the sequence \((f(x_n))_{n \geqslant 1}\) tends to \(f(x)\) . By possibly deleting the first terms of the sequence \((x_n)_{n \geqslant 1}\) , we reduce to the case where the terms of the sequence all belong to the same path-connected neighborhood of \(x\) in X. By the lemma below, there then exists a path \(c: \mathbf{I} \to \mathbf{X}\) such that \(c(0) = x\) and \(c(1/n) = x_n\) for \(n \geqslant 1\). If the map \(f\) is path-continuous, the map \(f \circ c\) is continuous and the element \(f(x) = f(c(0))\) is the limit of the sequence \((f(c(1/n)))_{n \geqslant 1}\) , that is, of the sequence \((f(x_n))_{n \geqslant 1}\), whence the corollary.

Lemma. --- Let X be a connected and locally path-connected topological space and let x be a point of X. Suppose that the point x has a countable fundamental system of neighborhoods. Then, for every sequence \((x_{n})_{n\geqslant1}\) of points of X tending to x, there exists a path c in X such that \(c(0)=x\) and \(c(1/n)=x_{n}\) for \(n\geqslant1\).

Let \((\mathrm{W}_{m})_{m\geqslant1}\) be a fundamental system of neighborhoods of x. Set \(V_{0}=X\) and for every \(m\geqslant1\), let \(V_{m}\) be a path-connected neighborhood of x contained in \(V_{m-1}\cap W_{m}\).

For every integer \(n \geqslant 1\), denote by \(m_n\) the largest integer \(m \leqslant n\) such that \(x_k \in V_m\) for all \(k \geqslant n\). The sequence \((m_n)_{n \geqslant 1}\) is increasing by definition; it tends to infinity, because the sequence \((x_n)_{n \geqslant 1}\) tends to \(x\). For every integer \(n \geqslant 1\), let \(c_n: \mathbf{I} \to V_{m_n}\) be a path of origin \(x_{n+1}\) and of endpoint \(x_n\) in \(V_{m_n}\). Let us define a map \(c: \mathbf{I} \to X\) by setting \(c(0) = x\) and \(c(t) = c_n(n(n+1)t - n)\) if \(1/(n+1) < t \leqslant 1/n\). We have \(c(1/n) = x_n\) for all \(n \geqslant 1\). The map \(c\) is therefore continuous on every interval of the form \([1/(n+1), 1/n]\) with \(n \geqslant 1\), hence on the interval \([0,1]\). If \(t \leqslant 1/n\), the point \(c(t)\) belongs to \(V_{m_n}\); the map \(c\) is therefore continuous at 0.

PROPOSITION 13. --- Let \(p \colon \mathrm{E} \to \mathrm{B}\) be a separated étale morphism and let \(s \colon \mathrm{B} \to \mathrm{E}\) be a section of \(p\). If the space B is locally path-connected and if the section \(s\) is path-continuous, it is continuous.

Let \(b\) be a point of B; let us prove that \(s\) is continuous at the point \(b\). Since \(p\) is an étale map, there exists a neighborhood V of \(b\) and a continuous local section \(s'\) of \(p\) defined in V such that \(s'(b) = s(b)\) (I, p. 33, prop. 9). Since B is locally path-connected, we may assume that V is path-connected. For every path \(c\) in V, originating from \(b\), the maps \(s \circ c\) and \(s' \circ c\) are two continuous liftings to E of the map \(c \colon \mathbf{I} \to \mathbf{B}\) and one has \(s \circ c(0) = s' \circ c(0) = s(b)\). By Corollary 1 of I, p. 34, we have \(s \circ c = s' \circ c\) and in particular \(s \circ c(1) = s' \circ c(1)\). Since every point of V is the endpoint of a path in V starting from \(b\), the maps \(s\) and \(s'\) coincide in V. The map \(s\) is therefore continuous in V.

COROLLARY. --- Let B be a topological space, and let (E, p) and (E', p') be two B-spaces. Suppose that the map p is étale and separated and that the space E' is locally path-connected. Then every path-continuous map \(f: E' \to E\) such that \(p \circ f = p'\) is continuous.

The map \(pr_{1}: E'\times_{B} E \rightarrow E'\) is étale and separated (I, p. 31, prop. 8 and I, p. 27, prop. 4), and the map \(x \mapsto (x, f(x))\) is a path-continuous section. By Proposition 13, it is continuous, hence f is continuous.

\subsection{Further results on compact metrizable topological spaces}

We endow the two-element set \(\{0,1\}\) with the discrete topology and the set \(\{0,1\}^{N}\) with the product topology. The topological space \(\{0,1\}^{N}\) is compact (TG, I, p. 63, th. 3), metrizable (TG, IX, p. 15, cor. 2), nonempty, totally disconnected (TG, I, p. 84, prop. 10), and has no isolated point.

PROPOSITION 14. --- Every compact, metrizable, nonempty, totally disconnected topological space with no isolated point is homeomorphic to \(\{0,1\}^{N}\) .

Let X be such a topological space. Endow it with a distance compatible with its topology. Since the space X is compact, it is complete for this distance (TG, II, p. 27, th. 1).

Lemma 3. --- Let \(\varepsilon\) be a real number \(>0\) . There exists an integer \(m \geqslant 1\) such that, for every integer \(n \geqslant m\) , X admits a partition consisting of \(n\) nonempty open and closed sets of diameter \(\leqslant \varepsilon\) .

Every point of X admits an open and closed neighborhood of diameter \(\leqslant\varepsilon\) (TG, II, p. 32, corollary of prop. 6). Since the space X is compact, it possesses a finite covering by such sets. Choose one of these, \((\mathrm{U}_{i})_{1\leqslant i\leqslant m}\) , for which m is minimal. We have \(m\geqslant1\) since X is not empty. For \(1\leqslant i\leqslant m\) , denote by \(V_{i}\) the intersection of \(U_{i}\) and the

\(X\setminus U_k\) for \(k < i\). Then \((V_i)_{1 \leqslant i \leqslant m}\) is a partition of X, consisting of m nonempty open and closed sets of diameter \(\leqslant \varepsilon\).

Since X has no isolated point, every nonempty open and closed subset V of X contains at least two points. Moreover, since X is compact and totally disconnected, V is the union of two disjoint nonempty open and closed subsets (loc. cit.). It follows, by induction, that for every integer \(n \geqslant m\) , X admits a partition consisting of n nonempty open and closed sets of diameter \(\leqslant \varepsilon\) .

Let us now finish the proof of Prop. 14. Every nonempty open and closed subspace of X is a compact, totally disconnected metric space with no isolated point. Lemma 3 therefore allows us to construct by induction a sequence \((\mathrm{J}_{n})_{n\geqslant0}\) of finite sets and, for every \(n\geqslant0\), a map \(\varphi_{n}\) from the set \(C_{n}=J_{0}\times\cdots\times J_{n}\) into the set of nonempty open and closed subsets of X of diameter \(\leqslant2^{-n}\), such that:

\begin{enumerate}
\def\labelenumi{(\roman{enumi})}
\item
  For every \(n \geqslant 0\) , there exists an integer \(m_n \geqslant 1\) such that \(\mathrm{Card}(\mathrm{J}_n) = 2^{m_n}\) ;
\item
  The family \((\varphi_0(c))_{c\in \mathrm{C}_0}\) is a partition of \(X\) ;
\item
  For every \(n \geqslant 0\) and every \(c \in \mathrm{C}_n\) , the family \((\varphi_{n+1}(c,j))_{j \in \mathrm{J}_{n+1}}\) is a partition of \(\varphi_n(c)\) .
\end{enumerate}

Denote by \(p_n\) the canonical projection of \(C_{n+1}\) onto \(C_n\). The sequence \(C = (C_n, p_n)_{n \geqslant 0}\) is a sieve (TG, IX, p. 63, def. 8). The topological space associated with this sieve (TG, IX, p. 63) is identified with the topological space J, the product of the discrete topological spaces \(J_n\). The sieve C and the sequence of maps \((\varphi_n)_{n \geqslant 0}\) define a strict sieving of the metric space X (TG, IX, p. 63 and p. 64). The map \(f: J \to X\) deduced from this sieving is continuous and bijective (TG, IX, p. 65). Since the topological space J is compact (TG, I, p. 63, th. 3) and X is Hausdorff, \(f\) is a homeomorphism (TG, I, p. 63, cor. 2). Since \(J_n\) is homeomorphic to \(\{0, 1\}^{m_n}\) for all \(n \geqslant 0\), J is homeomorphic to \(\{0, 1\}^N\) (TG, I, p. 25, prop. 2).

Example. --- Let K be the Cantor ternary set (TG, IV, p. 9, example). For every \(n \geqslant 0\) , set \(J_{n} = \{0,1\}\) and define a map \(K_{n}\) from the set \(C_{n} = J_{0} \times \cdots \times J_{n}\) into the set of closed intervals of {[}0,1{]} as follows: we set \(K_{0}(0) = [0, \frac{1}{3}]\) and \(K_{0}(1) = [\frac{2}{3}, 1]\) ; for every \(n \geqslant 0\) and every \(c \in C_{n}\) , \(K_{n+1}(c,0)\) and

\(K_{n+1}(c,1)\) are respectively the "left third" and the "right third" of \(K_n(c)\). If \(c = (j_0,j_1,\ldots,j_n) \in C_n\), \(K_n(c)\) is the interval denoted \(K_{n,p}\) in loc. cit., with \(p = 2^n j_0 + 2^{n-1}j_1 + \cdots + j_n + 1\), it is also the interval \([a,a + \frac{1}{3^{n+1}}]\), where \(a = 2(\frac{j_0}{3} + \frac{j_1}{3^2} + \cdots + \frac{j_n}{3^{n+1}})\). For \(n \geqslant 0\) and \(c \in C_n\), set \(\varphi_n(c) = K_n(c) \cap K\). The family \((\varphi_n(c))_{c \in C_n}\) is a partition of K formed by closed sets. These sets are therefore also open in K; they are nonempty and of diameter \(\frac{1}{3^{n+1}}\), because the endpoints of the intervals \(K_n(c)\) belong to K. The sequence \(C = (C_n, p_n)_{n \geqslant 0}\), where \(p_n: C_{n+1} \to C_n\) is the canonical projection \((c,j) \mapsto c\), is a sieve, and the topological space associated with this sieve is identified with \(\{0,1\}^N\). The sieve C and the sequence of maps \((\varphi_n)_{n \geqslant 0}\) define a strict sieving of the metric space K. The map \(f: \{0,1\}^N \to K\) deduced from this sieving is a homeomorphism, given by the formula

\[
f ((j _ {n}) _ {n \geqslant 0}) = 2 \sum_ {n = 0} ^ {\infty} \frac {j _ {n}}{3 ^ {n + 1}}.
\]

COROLLARY. --- Let X be a metrizable, compact, and nonempty topological space. There exists a continuous and surjective map from \(\{0,1\}^{N}\) to X.

Every compact and metrizable topological space is homeomorphic to a subspace, necessarily closed, of the topological space \(I^{N}\) (TG, IX, p. 18, prop. 12 and p. 21, prop. 16). Let A be a closed subspace of \(I^{N}\) and let h be a homeomorphism from A onto X.

Let us set \(K = \{0,1\}^{N}\) and, for \(\alpha = (a_n) \in K\), set \(f(\alpha) = \sum_{n=0}^{\infty} a_n 2^{-n-1}\); we have \(f(\alpha) \in I\). The map \(f: K \to I\) thus defined is surjective (TG, IV, p. 42) and continuous. Indeed, if two elements \(\alpha\) and \(\beta\) of K have the same coordinates of index \(< n\), we have \(|f(\beta) - f(\alpha)| \leqslant 2^{-n}\). Denote by \(g\) the map \((\alpha_n) \mapsto (f(\alpha_n))\) from \(K^N\) into \(I^N\); it is continuous and surjective. The topological space \(K^N\) is compact (TG, I, p. 63, th. 3), metrizable (TG, IX, p. 15, cor. 2), and totally disconnected (TG, I, p. 84, prop. 10), and the same holds for its closed subspace \(\bar{g}^{1}(A)\); the latter is nonempty since the maps \(g\) and \(h\) are surjective.

Then, the space \(\bar{g}^{1}(A) \times K\) is homeomorphic to \(\{0,1\}^{N}\) (prop. 14), since it is a compact, metrizable, totally disconnected topological space without isolated points, and the map \((x,y) \mapsto h(g(x))\) of \(\bar{g}^{1}(A) \times K\) into X is continuous and surjective.

\subsection{Topological properties of the image of a path}

PROPOSITION 15. --- The image of a path in a Hausdorff topological space is a compact, metrizable, connected, and locally path-connected topological space.

Let X be a Hausdorff topological space and \(c\colon I \to X\) a continuous map. Denote by R the equivalence relation \(c(s) = c(t)\) in I. The topological space \(c(\mathbf{I})\) is Hausdorff (TG, I, p. 63, cor. 1), the space I/R is quasi-compact (TG, I, p. 62, th. 2), hence the bijection I/R \(\to c(\mathbf{I})\) deduced from c is a homeomorphism (TG, I, p. 63, cor. 2). Consequently, the space \(c(\mathbf{I})\) is compact, metrizable (TG, IX, p. 22, prop. 17), connected (TG, I, p. 82, prop. 6), and locally path-connected (III, p. 261, prop. 8).

THEOREM 1 (Hahn and Mazurkiewicz). --- Every metrizable, compact, nonempty, connected, and locally connected topological space is homeomorphic to a quotient space of the interval {[}0,1{]}.

Lemma 4. --- Let K be a compact topological space and let R be a set of open sets of K covering K. There exists an entourage V of the uniform structure of K (TG, II, p. 27, th. 1) such that, for every \(x \in K\) , \(V(x)\) is contained in one of the sets belonging to R.

For every point x of K, there exists an entourage \(W_{x}\) of the uniform structure of K such that \(\mathrm{W}_{x}(x)\) is contained in one of the sets belonging to R. Let \(V_{x}\) be an entourage of the uniform structure of K such that \(\overset{2}{\mathrm{V}}_{x}\) is contained in \(W_{x}\). The interiors of the \(\mathrm{V}_{x}(x)\) cover K; since the space K is compact, there exists a finite subset F of K such that the family \((\mathrm{V}_{y}(y))_{y\in\mathrm{F}}\) is a covering of K (TG, I, p. 59). Denote by V the intersection of the family \((\mathrm{V}_{y})_{y\in\mathrm{F}}\); the set V is an entourage of the uniform structure of K. For every point \(x\in K\), there exists a point \(y\in F\) such that x belongs to \(\mathrm{V}_{y}(y)\). Consequently, the set \(\mathrm{V}(x)\) is contained in \(\overset{2}{\mathrm{V}}_{y}(y)\), hence in one of the sets belonging to R.

Lemma 5. --- Let X and Y be uniform spaces, and let f be a mapping from X into Y. Let F be a collection of closed subsets of X covering X and having the following properties:

\begin{enumerate}
\def\labelenumi{(\roman{enumi})}
\item
  There exists a set \(F_{0} \in F\) which meets all the sets \(F \in F\) ;
\item
  For every entourage U of the uniform structure of X, there are only finitely many sets \(F \in F\) which are not small of order U;
\item
  For every entourage V of the uniform structure of Y, there are only finitely many sets \(F \in F\) such that \(f(F)\) is not small of order V.
\end{enumerate}

Then, if the restriction of f to each of the sets \(F \in F\) is continuous, f is continuous.

Let x be a point of X. Let us prove that f is continuous at x. There exists a set \(F_{1} \in F\) such that \(x \in F_{1}\). The restriction of \(f\) to \(F_{0} \cup F_{1}\) is continuous (TG, I, p. 19, prop. 4). Replacing \(F_{0}\) by \(F_{0} \cup F_{1}\), we reduce to the case where \(x \in F_{0}\).

Let V be an entourage of the uniform structure of Y. Choose an entourage \(V'\) of this same uniform structure such that \(V' \subset V\). Since the restriction of f to \(F_0\) is continuous, there exists an entourage U of the uniform structure of X such that \(f(z) \in \mathrm{V}'(f(x))\) for all \(z \in \mathrm{F}_0 \cap \mathrm{U}(x)\). Let \(U'\) be an entourage of the uniform structure of X such that \(U' \subset U\). Let A denote the union of \(F_0\), of the sets \(F \in F\) which are not small of order \(U'\) and of those such that \(f(\mathrm{F})\) is not small of order \(V'\). By hypothesis, A is the union of a finite number of sets belonging to F, and the restriction of f to A is continuous (loc. cit.). There therefore exists a neighborhood W of x in X, contained in \(U'(x)\), such that \(f(y) \in \mathrm{V}(f(x))\) for \(y \in A \cap W\). To conclude, it will suffice to prove that we also have \(f(y) \in \mathrm{V}(f(x))\) for every point \(y \in (\mathrm{X} - \mathrm{A}) \cap \mathrm{W}\). Let y be such a point. Let F be an element of F such that \(y \in F\). By definition of A, F is small of order \(U'\) and \(f(\mathrm{F})\) is small of order \(V'\). By hypothesis, F meets \(F_0\). Let \(z \in F \cap F_0\). We have \(z \in U'(y)\) since F is small of order \(U'\) and \(y \in U'(x)\) since W is contained in \(U'(x)\), whence \(z \in U'(x)\) and a fortiori \(z \in U(x)\). But then, since z belongs to \(F_0\), we have \(f(z) \in V'(f(x))\). Moreover \(f(\mathrm{F})\) is small of order \(V'\), whence \(f(y) \in V'(f(z))\). It follows that one has \(f(y) \in V'(f(x))\) and finally \(f(y) \in V(f(x))\). This concludes the proof of Lemma 5.

Let us now prove Theorem 1. Let X be a nonempty compact, connected, and locally connected metric space. Such a space is path-connected and locally path-connected (III, p. 267, Corollary 2). By the corollary and the example of III, p. 270, there exists a continuous surjective map f from the Cantor ternary set K (TG, IV, p. 9, example) onto X. We shall construct a continuous extension g of f to {[}0,1{]}, which will prove Theorem 1.

Let \(\varepsilon\) be a real number \(>0\). The open, path-connected subsets of X of diameter \(\leqslant \varepsilon\) cover X. Let \(\mathcal{R}\) denote the set of their inverse images under \(f\); this is a set of open subsets of K covering K. By Lemma 4 of III, p. 272, there exists a real number \(\alpha > 0\) such that every closed ball in K of radius \(\alpha\) is contained in an element of \(\mathcal{R}\). In particular, if \(t\) and \(t'\) are points of K such that \(|t - t'| \leqslant \alpha\), there exists a path in X joining \(f(t)\) to \(f(t')\) whose image has diameter \(\leqslant \varepsilon\).

The preceding paragraph allows one to construct by induction a strictly increasing sequence \((n_{k})_{k\geqslant0}\) of integers \(\geqslant0\) having the following property: for every integer \(k\geqslant0\) and every pair \((t,t')\) of points of K such that \(|t-t'| \leqslant3^{-n_{k}}\), there exists a path in X joining \(f(t)\) to \(f(t')\) whose image has diameter \(\leqslant2^{-k}\). The complement of K in {[}0,1{]} is the union of a family \((\mathrm{I}_{n,p})\) of pairwise disjoint open intervals, where n ranges over the set of integers \(\geqslant0\) and p over the set of integers between 1 and \(2^{n}\) (TG, IV, p. 9, example). Consider one of these intervals \(I_{n,p}\) and write it as {]}a,b{[}. The points a and b belong to K and one has \(b-a=3^{-n-1}\). One defines the function g on the interval \(I_{n,p}\) in the following way: one chooses a path c in X joining \(f(a)\) to \(f(b)\), whose image has diameter \(\leqslant2^{-k}\) if \(n_{k}\leqslant n<n_{k+1}\), and one sets \(g(t)=c(\frac{t-a}{b-a})\) for \(t\in I_{n,p}\). The function \(g\colon[0,1]\to X\) thus defined extends f. It is continuous on K as well as on each of the closed intervals \(\overline{I_{n,p}}\). These latter meet K. Moreover, for every real number \(\varepsilon>0\), there are only finitely many intervals \(\overline{I_{n,p}}\) of length \(>\varepsilon\), and there are only finitely many intervals \(\overline{I_{n,p}}\) whose images under g have diameter \(>\varepsilon\). By Lemma 5, the map g is continuous.
